\chapter{Umsetzung}

In diesem Kapitel wird der Aufbau der DML Lautsprecher und die Montage im laborraum anhand der Messergebnisse aus Kapitel 3 Dokumentiert und dargestellt. Die DML Lautsprecher werden final an dieser Position hängen bleiben und in zukunft verwendung finden. 

\section{Bau nach Testergebnissen}

Durch die Auswertung der Versuchsergebnisse in Kapitel drei werden die Exciter auf den Platten nach Anhang \textbf{X} Angebracht. Als aufhängung der fäden werden Splinte verwendete welche mit einem Speziellen Kleber in den PUR Hartschaum eingeklebt werden. Die Kabel der Exciter werden an einen Hinter der Platte an der wand verklebten Kabelkanal befestigt und über die Hängedecke des Laborraums geführt. An der Hägngedecke werden die vom Lautsprecher kommenden Kabel mit den von der Stereoanlage kommenden lautsprecherkabeln verbunden. Dabei handelt es sich um \textbf{2,5mm2} Mehrlitzen kabel, die verbindung erfolgt über Litzenzugelassene Wagoklemmen mit zwei eingängen verbunden, die WAGO klemmen sind lutzen geeginet siehe Datenblatt Anhang\textbf{XY}. Die platte wird mit Angelschnüren an den Leisten der Hängedecke befestigt, als Befestigung dienen 3D Gedruckte hinter die Leisten greifende Harken aus PLA Plastik Anhang \textbf{XY}. Bei der Angelschnur handelt es sich um eine ... Da die Exciter unterhalb der Mittelline der Platte Positioniert werden entseht ein Kippmoment welches die Platte an der unterseite weiter von der wand weg drückt. Um diesem abstand entgegen zu wirken so das die Platte Lotrecht hängt werden an der den unteren ecken an der Hinteren seite der Platte 3D gedruckte Ösen Anhang\textbf{ XY} auf die Platte geklebt. diese zwei ösen werden mit der selben Angelschnur an eine dritte öse welche auf dem Kabelkanal verklebt ist verbunden womit die Platte gerade gezogen wird. 

\section{Subjektive Hör-Validierung}

zur subjektiven Hör validierung wurden verschiedene Musik Titel auf den DML Lautsrpechern wiedergegeben und von anhand 

\subsection{Ergebnis Diskussion}



